\documentclass{article}
\usepackage{amsmath,amssymb,amsthm}
\newtheorem{theorem}{Theorem}
\begin{document}
\title{Excluding the entire critical nine-degree-four profile by common holes}
\maketitle
\section{H01: Every original input gives the same intrinsic holes}
Let $G\subseteq Q_8$ be C4-free with683 edges, and let one specified
global parity$P$ have degree counts
\[
             (n_0,\ldots,n_8)=(0,0,0,0,9,71,45,2,1).       \tag{1}
\]
The opposite parity is unrestricted. Import the actual-face equality
theorem in q8\_equality\_profiles/main.tex, retaining its explicitly
declared historical316/317/321 local imports. The cost in(1) is1149,
so all112 actual coordinate$Q_6$ faces have$L_F=0$ or2 and
$2S_F+c_F=316$ or321, respectively. Here low means restricted degree
at most two on$P$; all unmarked roots have global degree at least five
and cannot be low after omitting two directions.

The nine degree-four roots contribute54 low-face incidences. There
are consequently27 faces with$L=2$ and85 with$L=0$.
For$L=2$, the odd weighted equality permits only$(c,S)=(1,160)$ or
$(3,159)$, since$c\le4$. The second violates the analytic theorem
$S+L\le158$ for exactly three full roots, with no degree premise on
either side (q6\_three\_full\_structure\_cycle010/PROOF.md).
Thus all27 double-low faces have zero specified$Q_3$ holes.
For$L=0$, the equality gives$c=0,2,4$ and exactly$2+c/2$ holes,
between two and four.

A specified hole is an actual coordinate$Q_3$ containing no full
specified three-star. C4-freeness makes the star-to-$Q_3$ map injective:
two distinct same-parity roots of a$Q_3$ have distance two, so two full
three-stars would retain a square. Thus holes are intrinsic to the
same original graph in every containing face. In the whole$Q_8$,
\[
 S_P=9\cdot4+71\cdot10+45\cdot20+2\cdot35+56=1772.
\]
There are1792 actual$Q_3$ faces, so$G$ has exactly20 specified holes.

\section{H02: A complete finite low-mark projection}
Use the necessary direction classification in
nine4\_direction\_structure/main.tex. After one common coordinate
permutation and root relabeling, the retained masks of the nine roots are
\[
       (170,204,102,240,90,60,150,T,T),\quad
                    T\in\{240,30,15,23\}.                  \tag{2}
\]
These are the seven nonzero linear-function1-fibres on$\mathbb F_2^3$
and the four previously proved types of a four-subset$T$.
Let$v_j$ be the actual position of root$j$. Its low faces depend only
on the36 retained-coordinate bits
\[
                   b_{ja}=(v_j)_a,\qquad a\in R_j.         \tag{3}
\]
Indeed a face with omitted pair$I$ and fixed bits$z$ has this root low
exactly when$I\subseteq R_j$ and$(v_j)_I=z$. The other position bits
play no role in deciding its low faces. We drop their constraints,
root distinctness and common parity in this projection; admitting extra
marks is safe for a necessary condition.

For each coordinate$a$ occurring in any mask, let$j(a)$ be the first
such root in(2). One simultaneous translation by the bits
$d_a=b_{j(a),a}$ makes all$b_{j(a),a}$ zero. Unused coordinates may
have$d_a=0$. This is one actual cube translation, not a separate move
of each root. It transports every face and hole together and may
exchange the two parity names, which does not change the argument.

For an omitted pair$\{a,b\}$, the number$m$ of masks containing it
is0,2 or4. Even multiplicities in each fixed-bit cell imply
\[
       \bigoplus_{j:\{a,b\}\subseteq R_j} b_{ja}=0,
       \qquad
       \bigoplus_{j:\{a,b\}\subseteq R_j} b_{jb}=0.        \tag{4}
\]
Together with the common translation gauge these are homogeneous
linear equations on all36 bits. Ordinary exact binary elimination
gives dimensions8,4,3,3 for the four values of$T$ in(2).
Enumerate all$256,16,8,8$ solutions, respectively, then count every
one of the112 actual cells and keep exactly those with count0 or2.
The resulting numbers of retained-coordinate mark cases are15,5,4,4.
This complete28-case cover, inducing22 distinct sets of double-low faces,
is saved in ALL\_LOW\_MARKS.json and produced
by enumerate\_low\_marks.py. No independent realizability of its marks
is asserted or needed. Every original graph gives one enumerated
pattern after the common transformations above.

\section{H03: Common-hole requirements and elementary rejections}
For each enumerated pattern let$\mathcal B$ be its27 actual double-low
faces, and$\mathcal U$ all actual$Q_3$ faces contained in none of them.
By H01 every original hole belongs to$\mathcal U$. There must be one
set$X\subseteq\mathcal U$ with
\[
  |X|=20,\qquad 2\le|\{Q\in X:Q\subseteq F\}|\le4
                \quad\text{for all85 zero-low faces }F.    \tag{5}
\]
Only literal face inclusion is used. No unknown ordinary star or
separately feasible face completion is chosen. We do not even exclude
holes filled by the nine known stars; dropping that further necessary
condition enlarges the domain in(5).

For18 enumerated mark cases a zero-low face has fewer than two elements of
$\mathcal U$, immediately violating(5). For four more mark cases,
every zero-low face with exactly two available elements forces both
into$X$. Their union consists of24 distinct$Q_3$ faces, contradicting
$|X|=20$. The complete lists and witnesses are in
ALL\_PATTERN\_RESULTS.json. This leaves three mark cases for each of
$T=15,23$, six in total. Their$\mathcal U$ have252 elements.

\section{H04: Exact rejection of the remaining six common domains}
For each remaining mark case introduce one Boolean variable$x_Q$ for
every$Q\in\mathcal U$. Encode(5) by exact prefix threshold gates,
as in hole\_layout.py and all\_mark\_holes.py. The threshold
$u_{i,k}$ means at least$k$ of the first$i$ inputs are true:
\[
 u_{i,1}=u_{i-1,1}\lor x_i,\qquad
 u_{i,k}=u_{i-1,k}\lor(u_{i-1,k-1}\land x_i),
\]
with missing thresholds initially false and the obvious first-input
boundary. Each AND/OR gate is encoded in both directions. For a
range$[\ell,h]$ retain thresholds through$h+1$, asserting$u_{n,\ell}$
when$\ell>0$ and$\neg u_{n,h+1}$ when$h<n$. Thus an input assignment
extends to the gate clauses if and only if its count is in that range;
induction on$i$ supplies the proof and every auxiliary value.
The global range is$[20,20]$; each actual good-face range is$[2,4]$.
Complete short-counter truth calibration covers7423 assignments.

The six CNFs were generated and dispatched once each. All returned
UNSAT, and every trace was strictly replayed against its own frozen
CNF with the existing check\_rup\_cached program. The six certified
addition counts are
\[
                   8963,6970,7919,6754,9369,9944,
\]
total49919. Every replay validates a final empty clause. Per-case
INPUT.json lists every actual$Q_3$ variable and every face incidence;
STARTED.json and RESULT.json bind CNF, solver/checker and trace hashes.
Replay establishes these finite CNFs inconsistent; H02 and the exact
encoding bridge, rather than solver failure, connect them to all
remaining original inputs.

\section{H05: Return to the original graph and limits}
\begin{theorem}
No683-edge C4-free subgraph of$Q_8$ has profile(1) on either specified
parity; the opposite parity is unrestricted.
\end{theorem}
\begin{proof}
Any such actual graph supplies the one intrinsic20-hole set of H01.
The common direction/translation reduction places its low faces in
one of the28 exhaustive mark cases of H02. Eighteen and four cases
are impossible by H03; the remaining six exact common-hole domains
are inconsistent by H04. This exhausts every input and is a
contradiction.
\end{proof}
The earlier36-edge partial marked patterns remain valid partial graphs.
Their separate per-face availability did not ensure a common20-hole
set. This theorem does not retract those projection counterexamples or
turn any partial graph into a683 graph. It strictly strengthens that
test by using a single shared hole set.

This is a candidate local computational proof requiring fresh review
of the complete cover, actual input bridge, encoder and replays.
Historical local imports are retained rather than replayed. No Lean,
TeX/PDF compilation or full unrestricted682 proof is asserted here.
\end{document}
